\documentclass{article}
\usepackage[
    mocha,
    pagecolor=true,
    textcolor=true
]{catppuccinpalette}
\usepackage[
    backend=biber,
    style=alphabetic,
    sorting=ynt
]{biblatex}
\usepackage{mdframed}
\usepackage{letltxmacro}
\usepackage{xcolor}
\usepackage{enumitem}
\usepackage[fixed]{fontawesome5}
\usepackage{titlesec}
\usepackage{abstract}
\usepackage{hyperref}
\usepackage[nottoc,numbib]{tocbibind}
\usepackage{tikz}
\usepackage{pgfplots}
\usepackage{caption}
\usepackage{multicol}
\usepackage{amsmath}

% Pgfplots 
\pgfplotsset{compat=1.18}

% Catppuccin mocha colors
\definecolor{rosewater}{HTML}{f5e0dc}   % Rosewater
\definecolor{flamingo}{HTML}{f2cdcd}    % Flamingo
\definecolor{pink}{HTML}{f5c2e7}        % Pink
\definecolor{mauve}{HTML}{cba6f7}       % Mauve
\definecolor{red}{HTML}{f38ba8}         % Red
\definecolor{maroon}{HTML}{eba0ac}      % Maroon
\definecolor{peach}{HTML}{fab387}       % Peach
\definecolor{yellow}{HTML}{f9e2af}      % Yellow
\definecolor{green}{HTML}{a6e3a1}       % Green
\definecolor{teal}{HTML}{94e2d5}        % Teal
\definecolor{sky}{HTML}{89dceb}         % Sky
\definecolor{sapphire}{HTML}{74c7ec}    % Sapphire
\definecolor{blue}{HTML}{89b4fa}        % Blue
\definecolor{lavender}{HTML}{b4befe}    % Lavender
\definecolor{text}{HTML}{cdd6f4}        % Text
\definecolor{subtext1}{HTML}{bac2de}    % Subtext 1
\definecolor{subtext0}{HTML}{a6adc8}    % Subtext 0
\definecolor{overlay2}{HTML}{9399b2}    % Overlay 2
\definecolor{overlay1}{HTML}{7f849c}    % Overlay 1
\definecolor{overlay0}{HTML}{6c7086}    % Overlay 0
\definecolor{surface2}{HTML}{585b70}    % Surface 2
\definecolor{surface1}{HTML}{45475a}    % Surface 1
\definecolor{surface0}{HTML}{313244}    % Surface 0
\definecolor{base}{HTML}{1e1e2e}        % Base
\definecolor{mantle}{HTML}{181825}      % Mantle
\definecolor{crust}{HTML}{11111b}       % Crust

% HREF
\hypersetup{
    colorlinks=true,
    linkcolor=peach,
    filecolor=yellow,
    urlcolor=sapphire
}

\urlstyle{same}

% Formatting
\titleformat{\section}
    {\normalfont\Large\bfseries\color{maroon}}
    {\thesection}{1em}{}

\titleformat{\subsection}
    {\normalfont\large\bfseries\color{peach}}
    {\thesubsection}{1em}{}

\titleformat{\subsubsection}
    {\normalfont\normalsize\bfseries\color{yellow}}
    {\thesubsubsection}{1em}{}

\renewcommand{\abstractname}{\color{red}\normalfont\normalsize\bfseries Abstract}

\setlist[enumerate,1]{label=\color{peach}\arabic*., font=\bfseries}
\setlist[enumerate,2]{label=\color{yellow}(\alph*), font=\bfseries}
\setlist[enumerate,3]{label=\color{green}\roman*., font=\bfseries}
\setlist[enumerate,4]{label=\color{teal}\Alph*., font=\bfseries}

\LetLtxMacro{\oldtextbf}{\textbf}
\renewcommand{\textbf}[1]{\oldtextbf{\color{mauve}#1}}
\LetLtxMacro{\oldtextit}{\textit}
\renewcommand{\textit}[1]{\oldtextit{\color{flamingo}#1}}

\everymath{\color{pink}}
\everydisplay{\color{pink}}

% Tikz Axis
\pgfplotsset{
    custom/.style={
        enlargelimits = true,
        thick,
        font = \small,
        axis line style = {rosewater},
        tick style = {rosewater},
        label style = {color = rosewater},
        xticklabel style = {color = rosewater},
        yticklabel style = {color = rosewater},
        legend pos = north east,
        legend style = {
            font=\small,
            draw=rosewater,
            fill=mantle
        }
    }
}

% Callouts
\mdfdefinestyle{calloutbase}{
  topline=false, rightline=false, bottomline=false,
  linewidth=2pt,
  innerleftmargin=10pt,
  innerrightmargin=10pt,
  innertopmargin=6pt,
  innerbottommargin=6pt,
  backgroundcolor=mantle,
  font=\color{subtext1},
}

\newmdenv[style=calloutbase, linecolor=sapphire]{cnote-base}
\newenvironment{cnote}[1][Note]{%
  \begin{cnote-base}%
  {\color{sapphire}\bfseries\faInfoCircle\hspace{0.5em}#1}\par\vspace{4pt}%
}{%
  \end{cnote-base}%
}

\newmdenv[style=calloutbase, linecolor=green]{ctip-base}
\newenvironment{ctip}[1][Tip]{%
  \begin{ctip-base}%
  {\color{green}\bfseries\faLightbulb\hspace{0.5em}#1}\par\vspace{4pt}%
}{%
  \end{ctip-base}%
}

\newmdenv[style=calloutbase, linecolor=peach]{cwarning-base}
\newenvironment{cwarning}[1][Warning]{%
  \begin{cwarning-base}%
  {\color{peach}\bfseries\faExclamationTriangle\hspace{0.5em}#1}\par\vspace{4pt}%
}{%
  \end{cwarning-base}%
}

\newmdenv[style=calloutbase, linecolor=overlay2]{cinfo-base}
\newenvironment{cinfo}[1][Further Information]{%
  \begin{cinfo-base}%
  {\color{overlay2}\bfseries\faBookOpen\hspace{0.5em}#1}\par\vspace{4pt}%
}{%
  \end{cinfo-base}%
}


\begin{document}
\title{Audió Szintetizátor Specifikáció}
\author{Pászti Tibor Miklós}
\maketitle

\section{Specifikáció:}
\subsection{Felhasználói Felület:}
Az Audió Szintetizátor alapvetően egy SDL3 GUI-n keresztül lesz 
elérhető, ahol az alábbi műveleteket lehet elvégezni:

\begin{enumerate}
    \item {MIDI fájl kiválasztás} 
    \item {Oszcillátor konfiguráció kiválasztás} 
    \item {Oszcillátor paraméterek módosítása (\textbf{waveform}, \textbf{velocity}, ...)} 
    \item {Audió fájl lejátszása} 
\end{enumerate}

\begin{cnote}[frametitle=Megjegyzés:]
    SDL3 főleg a megújult audió lejátszás miatt lett választva.
\end{cnote}

\begin{cnote}[frametitle=Megjegyzés:]
    A projekt szigorúan csakis a 
    \href{https://catppuccin.com/palette/}{Catppuccin Mocha}
    szín palettát fogja használni.
\end{cnote}

\subsection{MIDI Fájl:}
A program egy adott MIDI fájl alapján fogja le szintetizálni 
az audió fájl-t.

\begin{cnote}[frametitle=Megjegyzés:]
    Vannak MIDI eventek amikre nem lesz szükség a jelenlegi
    projekt keretein belül. 
    Erre néhány példa:
    \begin{enumerate}
        \item{Control/Mode Change}
        \item{Program Change}
        \item{Pitch Bend Change}
    \end{enumerate}

    \href{https://midi.org/expanded-midi-1-0-messages-list}{MIDI eventek}
\end{cnote}

\subsection{Oszcillátor:}
Minden MIDI channel kap egy saját oszcillátort aminek az 
alap állása egy Sine waveform, konstans velocity-vel.
A GUI-n keresztül lehetőség lesz az adott oszcillátor
paramétereinek módosítására, a jelenlegi paraméterek 
elmentésére és az addig elmentett konfigurációk beolvasására.

\subsubsection{Konfigurációs Fájl:}
Az oszcillátorok paramétereinek tárolása 
\href{https://en.wikipedia.org/wiki/INI_file}{INI}
fájlok-al fog történni, a könnyű olvashatóság és módosítás illetve kezelhető beolvasás
érdekében.

\subsubsection{Paraméterek:}
\begin{enumerate}
    \item{Waveform}
    \item{Velocity}
    \item{Amplitude}
\end{enumerate}

\subsubsection{Waveform-ok:}
Alapvetően csak 3 hullám tipus lesz:
\begin{enumerate}
    \item{Sine wave}
    \item{Square wave}
    \item{Triangle wave}
\end{enumerate}
\end{document}
